\documentclass[11pt]{article}
\usepackage[margin=1in]{geometry}
\usepackage{amsmath,amssymb,mathtools,bm}

\newcommand{\Exp}{\operatorname{Exp}}
\newcommand{\Log}{\operatorname{Log}}
\newcommand{\SO}{\mathrm{SO}}
\newcommand{\Ad}{\mathbf{Ad}}
\newcommand{\R}{\mathbb{R}}


\begin{document}

\section*{Problem}

Let
\[
R_1,R_2\in\SO(3).
\]
Calculate
\[
\boxed{
\frac{\partial(R_1R_2^{-1})}{\partial R_2}
}
\]
using both left and right perturbation models.

This note first derives the result for the \emph{right perturbation model}.
The left-perturbation derivation will be added separately.

\section{Right perturbation}

\subsection{Method 1: Directly from the definition}

Use the right-plus convention
\[
R_2' = R_2\Exp(\delta\theta),
\qquad
\delta\theta\in\R^3.
\]

Define
\[
Y(R_2)=R_1R_2^{-1}.
\]

Since the output $Y\in\SO(3)$ is also a Lie-group element, its
difference is measured with the right-minus operator
\[
Y'\ominus Y=\Log(Y^{-1}Y').
\]

To write the Jacobian limit rigorously, choose an arbitrary direction
\[
u\in\R^3
\]
and set
\[
\delta\theta=hu,
\qquad h\to0.
\]

The right Jacobian is therefore defined by
\[
\boxed{
{}^{R_2}J_{R_2}^{R_1R_2^{-1}}u
=
\lim_{h\to0}
\frac{
Y(R_2\Exp(hu))\ominus Y(R_2)
}{h}.
}
\tag{1}
\]

Substituting
\[
Y(R_2)=R_1R_2^{-1},
\]
gives
\[
{}^{R_2}J_{R_2}^{R_1R_2^{-1}}u
=
\lim_{h\to0}
\frac{
\Log\!\left[
(R_1R_2^{-1})^{-1}
R_1(R_2\Exp(hu))^{-1}
\right]
}{h}.
\tag{2}
\]

Using
\[
(R_1R_2^{-1})^{-1}=R_2R_1^{-1}
\]
and
\[
(R_2\Exp(hu))^{-1}
=
\Exp(-hu)R_2^{-1},
\]
the argument of the logarithm becomes
\begin{align}
&(R_1R_2^{-1})^{-1}
R_1(R_2\Exp(hu))^{-1}
\\
&=
R_2R_1^{-1}
R_1
\Exp(-hu)R_2^{-1}
\\
&=
R_2\Exp(-hu)R_2^{-1}.
\end{align}

Hence
\[
{}^{R_2}J_{R_2}^{R_1R_2^{-1}}u
=
\lim_{h\to0}
\frac{
\Log\!\left(
R_2\Exp(-hu)R_2^{-1}
\right)
}{h}.
\tag{3}
\]

Now use the conjugation identity
\[
\boxed{
R\Exp(a)R^{-1}
=
\Exp(\Ad_R a).
}
\tag{4}
\]

Therefore
\[
R_2\Exp(-hu)R_2^{-1}
=
\Exp\!\left(-h\Ad_{R_2}u\right).
\]

Substituting into (3),
\[
{}^{R_2}J_{R_2}^{R_1R_2^{-1}}u
=
\lim_{h\to0}
\frac{
\Log\!\left[
\Exp(-h\Ad_{R_2}u)
\right]
}{h}.
\]

Near the identity,
\[
\Log(\Exp(v))=v,
\]
so
\[
{}^{R_2}J_{R_2}^{R_1R_2^{-1}}u
=
-\Ad_{R_2}u.
\]

Since $u$ is arbitrary,
\[
\boxed{
{}^{R_2}J_{R_2}^{R_1R_2^{-1}}
=
-\Ad_{R_2}.
}
\tag{5}
\]

For $\SO(3)$,
\[
\Ad_R=R.
\]

Therefore the final right-perturbation result is
\[
\boxed{
{}^{R_2}J_{R_2}^{R_1R_2^{-1}}
=
-R_2.
}
\tag{6}
\]


\section{Left perturbation}

Use the left-plus convention
\[
R_2'=\Exp(\delta\theta)R_2.
\]

\subsection{Method 1: From equation (57)}

From the right-perturbation result,
\[
\boxed{
{}^{R_2}J_{R_2}^{R_1R_2^{-1}}
=
-R_2.
}
\]

Equation (57) converts a right Jacobian into the corresponding left
Jacobian:
\[
\boxed{
{}^{\varepsilon}J_X^{f(X)}
=
\Ad_{f(X)}
\;{}^{X}J_X^{f(X)}
\;\Ad_X^{-1}.
}
\tag{57}
\]

For
\[
f(R_2)=R_1R_2^{-1},
\]
this gives
\[
{}^{\varepsilon}J_{R_2}^{R_1R_2^{-1}}
=
\Ad_{R_1R_2^{-1}}
\left(-R_2\right)
\Ad_{R_2}^{-1}.
\]

For $\SO(3)$,
\[
\Ad_R=R,
\qquad
\Ad_R^{-1}=R^{-1}.
\]

Therefore
\begin{align}
{}^{\varepsilon}J_{R_2}^{R_1R_2^{-1}}
&=
(R_1R_2^{-1})(-R_2)R_2^{-1}
\\
&=
\boxed{-R_1R_2^{-1}}.
\end{align}

\subsection{Method 2: Directly from the definition}

Choose an arbitrary direction
\[
u\in\R^3
\]
and let
\[
\delta\theta=hu,
\qquad h\to0.
\]

Define
\[
Y=R_1R_2^{-1}.
\]

Under the left perturbation,
\[
R_2'=\Exp(hu)R_2.
\]

Hence
\[
Y'
=
R_1(R_2')^{-1}
=
R_1(\Exp(hu)R_2)^{-1}.
\]

Using
\[
(AB)^{-1}=B^{-1}A^{-1},
\]
we obtain
\[
(\Exp(hu)R_2)^{-1}
=
R_2^{-1}\Exp(-hu),
\]
so
\[
\boxed{
Y'
=
R_1R_2^{-1}\Exp(-hu)
=
Y\Exp(-hu).
}
\]

Because the output perturbation is left/global, use the left-minus
operator:
\[
Y'\ominus_LY
=
\Log(Y'Y^{-1}).
\]

Therefore
\begin{align}
{}^{\varepsilon}J_{R_2}^{R_1R_2^{-1}}u
&=
\lim_{h\to0}
\frac{
\Log(Y'Y^{-1})
}{h}
\\
&=
\lim_{h\to0}
\frac{
\Log\!\left(
Y\Exp(-hu)Y^{-1}
\right)
}{h}.
\end{align}

Using the conjugation identity
\[
Y\Exp(v)Y^{-1}
=
\Exp(\Ad_Yv),
\]
we have
\[
Y\Exp(-hu)Y^{-1}
=
\Exp(-h\Ad_Yu).
\]

Thus
\begin{align}
{}^{\varepsilon}J_{R_2}^{R_1R_2^{-1}}u
&=
\lim_{h\to0}
\frac{
\Log\!\left(
\Exp(-h\Ad_Yu)
\right)
}{h}
\\
&=
-\Ad_Yu.
\end{align}

Since $u$ is arbitrary,
\[
{}^{\varepsilon}J_{R_2}^{R_1R_2^{-1}}
=
-\Ad_Y.
\]

Finally, because
\[
Y=R_1R_2^{-1}
\]
and $\Ad_Y=Y$ for $\SO(3)$,
\[
\boxed{
{}^{\varepsilon}J_{R_2}^{R_1R_2^{-1}}
=
-R_1R_2^{-1}.
}
\]

\section*{Final comparison}

For the right perturbation
\[
R_2'=R_2\Exp(\delta\theta_R),
\]
\[
\boxed{
{}^{R_2}J_{R_2}^{R_1R_2^{-1}}
=
-R_2.
}
\]

For the left perturbation
\[
R_2'=\Exp(\delta\theta_L)R_2,
\]
\[
\boxed{
{}^{\varepsilon}J_{R_2}^{R_1R_2^{-1}}
=
-R_1R_2^{-1}.
}
\]

\end{document}
